\documentclass[11pt]{article}
\usepackage[margin=2.5cm]{geometry}
\usepackage{amsmath,amssymb,amsthm}
\usepackage{mathpartir}
\usepackage{xcolor}

\newenvironment{theorem}[1]{\par\medskip\noindent\textbf{Theorem #1.}\ \itshape}{\par\medskip}
\newenvironment{lemma}[1]{\par\medskip\noindent\textbf{Lemma #1.}\ \itshape}{\par\medskip}
\theoremstyle{definition}
\newtheorem{definition}{Definition}
\theoremstyle{remark}
\newtheorem*{remark}{Remark}

\newcommand{\vdashG}{\vdash_{\mathsf G}}
\newcommand{\vdashT}{\vdash_{\mathsf T}}
\newcommand{\U}{\mathsf U}
\newcommand{\T}{\mathsf T}
\newcommand{\code}[1]{\ulcorner #1\urcorner}
\newcommand{\p}{\mathsf p}
\newcommand{\q}{\mathsf q}
\newcommand{\id}{\mathsf{id}}
\newcommand{\lift}{\mathsf{lift}}
\newcommand{\sect}{\mathsf{sec}}
\newcommand{\app}{\mathsf{app}}
\newcommand{\er}[1]{|#1|_{\mathsf T}}
\newcommand{\ev}[2]{[\![#1]\!]#2}
\newcommand{\up}{\mathsf{up}}
\newcommand{\idv}{\iota}
\newcommand{\seg}[1]{\texttt{#1}}
\newcommand{\vdashR}{\vdash_{\mathsf R}}
\newcommand{\erR}[1]{|#1|_{\mathsf R}}

\title{Explicit and implicit substitutions with Ehrhard's generators\\[2pt]
\large A variant of Section 4 of ``From Tarski to Russell'', and Section 5, formalised in Unit}
\author{}
\date{22 November 2026}

\begin{document}
\maketitle

\noindent
Section 4 of Coquand--Sterbac proves that $T_{\mathsf{GAT}}$, the fully
explicit generalised algebraic theory with Dybjer's substitutions
$\id,\ \sigma\delta,\ (),\ (\sigma,a),\ \p,\ \q$, is equivalent to
$T_{\mathsf{Tarski}}$, where substitution is an operation on de Bruijn
terms. We replace $T_{\mathsf{GAT}}$ by a theory $T_{\mathsf E}$ whose
substitutions are generated as in Ehrhard's presentation of categories with
families: $\id,\ \sigma\delta,\ (),\ \p,\ \lift_A(\sigma),\ \sect_A(a)$.
We prove the analogues of Lemmas 15--19 and Theorem~20. We then prove
Section~5 of the paper, the equivalence of $T_{\mathsf{Tarski}}$ and
$T_{\mathsf{Russell}}$: Lemma~21, Lemmas~25--26, Corollary~27 and
Theorem~28 (Sections~\ref{sec:inv}--\ref{sec:section}), and compose them
with Theorem~20 to relate $T_{\mathsf{Russell}}$ and $T_{\mathsf E}$
directly (Section~\ref{sec:compose}). Finally, we use this relation to
interpret $T_{\mathsf{Russell}}$ in Ehrhard models, including the
iterative-set models (Section~\ref{sec:rsem}). The formal development
is checked in Unit by \seg{UNIT/GATCheck.seg}.

\section{The theory $T_{\mathsf E}$}

Raw expressions (contexts, types, terms and substitutions share one syntax):
\begin{align*}
\Gamma,\Delta &::= 1 \mid \Gamma.A \\
A,B &::= \Pi(A,B) \mid \U_n \mid \T_n(a) \mid A\sigma\\
a,b,c &::= \lambda(A,B,b) \mid \app(A,B,c,a) \mid \Pi_n(a,b) \mid \U^n_m \mid t^n_m(a) \mid \q \mid a\sigma \\
\sigma,\delta &::= \id \mid \sigma\delta \mid () \mid \p \mid \lift_A(\sigma) \mid \sect_A(a)
\end{align*}
The judgements are $\vdashG\Gamma$, $\vdashG\Gamma=\Delta$,
$\Gamma\vdashG A$, $\Gamma\vdashG a:A$, $\vdashG\sigma:\Delta\to\Gamma$ and
the corresponding equalities. The substitution rules are
\begin{mathpar}
\inferrule{\vdashG\Gamma}{\vdashG\id:\Gamma\to\Gamma}\and
\inferrule{\Gamma\vdashG A}{\vdashG\p:\Gamma.A\to\Gamma}\and
\inferrule{\Gamma\vdashG A}{\Gamma.A\vdashG\q:A\p}\and
\inferrule{\vdashG\sigma:\Delta\to\Gamma \\ \Gamma\vdashG A}{\vdashG\lift_A(\sigma):\Delta.A\sigma\to\Gamma.A}\and
\inferrule{\Gamma\vdashG A \\ \Gamma\vdashG a:A}{\vdashG\sect_A(a):\Gamma\to\Gamma.A}
\end{mathpar}
and the type of an application is $\Gamma\vdashG\app(A,B,c,a):B\,\sect_A(a)$.
Besides the category laws and the uniqueness of $()$, the equations are
Ehrhard's:
\begin{align*}
&\lift_A(\id)=\id, &&\lift_A(\sigma\delta)=\lift_A(\sigma)\,\lift_{A\sigma}(\delta),
&&\p\,\lift_A(\sigma)=\sigma\,\p, &&\p\,\sect_A(a)=\id,\\
&\sect_A(a)\,\sigma=\lift_A(\sigma)\,\sect_{A\sigma}(a\sigma),
&&\lift_A(\p)\,\sect_{A\p}(\q)=\id, &&\q\,\lift_A(\sigma)=\q, &&\q\,\sect_A(a)=a.
\end{align*}
The naturality equations are those of the paper with $\sigma^+$ replaced by
$\lift_A(\sigma)$, e.g.\ $\Pi(A,B)\sigma=\Pi(A\sigma,B\,\lift_A(\sigma))$
and $\app(A,B,c,a)\sigma=\app(A\sigma,B\lift_A(\sigma),c\sigma,a\sigma)$.
The other rules are $\beta$, $\eta$ in the form
$c=\lambda(A,B,\app(A\p,B\lift_A(\p),c\p,\q))$, and the six Tarski
equations. Finally there are congruence rules, and rules converting
judgements along $\vdashG\Gamma=\Delta$, including both ends of a
substitution:
\[
\inferrule{\vdashG\sigma:\Delta\to\Gamma \\ \vdashG\Delta=\Delta' \\ \vdashG\Gamma=\Gamma' \\ \vdashG\Delta' \\ \vdashG\Gamma'}
  {\vdashG\sigma:\Delta'\to\Gamma'}
\]
This last rule is needed for the presuppositions (Section~\ref{sec:presup}):
without it the right-hand side of $\lift_A(\id)=\id:\Gamma.A\id\to\Gamma.A$
is not typable, $\id$ having only the types $\Gamma\to\Gamma$.
Dybjer's pairing is definable: $(\sigma,a):=\lift_A(\sigma)\sect_{A\sigma}(a)$.

\begin{remark}
Raw terms carry no context annotation: $\U_n$, $\U^n_m$, $\id$, $()$, $\p$
and $\q$ are bare, and the context of the judgement serves as the annotation.
The type annotations $\lambda(A,B,b)$, $\app(A,B,c,a)$, $\lift_A$, $\sect_A$
are kept. This is what makes the lifting below independent of the context.
With annotated $\p_{\Gamma,A}$, $\q_{\Gamma,A}$, the relational form of
Lemma~18 would need presuppositions about the right-hand side of each
equation.
\end{remark}

\section{Erasure (Theorem 16)}

$T_{\mathsf{Tarski}}$ is the system of Appendix~A.2 over unscoped de Bruijn
terms; a substitution is a map $\alpha:\mathbb N\to\mathsf{Tm}$. Erasure
forgets the explicit substitutions:
\begin{gather*}
\er{A\sigma}=\er A[\er\sigma],\qquad \er{a\sigma}=\er a[\er\sigma],\qquad \er\q=v_0,\qquad
\er\id=\er{()}=\mathrm{id},\qquad \er\p=\mathrm{wk},\\
\er{\sigma\delta}=\er\sigma\circ\er\delta,\qquad
\er{\lift_A(\sigma)}=\er\sigma^+,\qquad \er{\sect_A(a)}=[\er a].
\end{gather*}
So each generator erases to one of the operations that the typing rules
of $T_{\mathsf{Tarski}}$ already use: weakening, lifting under a binder,
and single substitution.

\begin{theorem}{16}
If $\Gamma\vdashG J$ then $\er\Gamma\vdashT\er J$. Typed substitutions go to
typed substitutions, and equal ones to pointwise convertible ones.
\end{theorem}
The proof uses Lemma~15, the substitution lemma, in both its parts. Typed
substitutions $\Delta\vdashT\alpha:\Gamma$ are defined by recursion on
$\Gamma$, following the paper's variable rules. A single typed traversal over
a ``kit'' gives renaming, which needs only the variable rules, and then
substitution. Every Ehrhard equation erases to a pointwise syntactic
identity. For instance, $\er{\sect_A(a)\sigma}=\er{\lift_A(\sigma)\sect_{A\sigma}(a\sigma)}$
is the identity $([a])[\alpha]=\alpha^+\circ[a[\alpha]]$.

\section{Lifting by evaluation (Lemma 17)}

Instead of $\lceil-\rceil_\Gamma$ we evaluate a raw term of
$T_{\mathsf{Tarski}}$ in an \emph{environment} $\rho:\mathbb N\to T_{\mathsf E}$:
\[
\ev{v_i}\rho=\rho(i),\qquad
\ev{\Pi(A,B)}\rho=\Pi(\ev A\rho,\ev B{(\up\rho)}),\qquad
\up\rho=(\q,\ \rho(0)\p,\ \rho(1)\p,\dots),
\]
and homomorphically otherwise. Set
$\lceil t\rceil:=\ev t\idv$ with $\idv(i)=\q\,\p^i$.
Since the environment under a binder is $\up\rho$, the paper's identity
$(\lceil\alpha\rceil)^+=\lceil\alpha^+\rceil$ never has to be proved; it is
built into the definition.

\begin{lemma}{17}
$\er{\ev t\rho}=t[\er\rho]$, hence $\er{\lceil t\rceil}=t$ and
$\er{\lceil\Gamma\rceil}=\Gamma$, strictly.
\end{lemma}
Evaluation commutes strictly with renamings,
$\ev{t[r]}\rho=\ev t{(\rho\circ r)}$; in particular
$\ev{t[\mathrm{wk}]}\rho=\ev t{(\rho\circ\mathrm{suc})}$.

\section{Lemma 18}

An environment $\rho$ is \emph{typed} over $\Delta$ for $\Gamma$ when
$\Delta\vdashG\rho(i):\ev{A_i}\rho_{>i}$ for each type $A_i$ of $\Gamma$.
Two environments are \emph{related} when, in addition, $\rho_1(i)=\rho_2(i)$
pointwise. For $\Gamma\vdashT A$ consider four properties: for all typed $\rho$,
\begin{itemize}
\item[(P)] $\Delta\vdashG\ev A\rho$;
\item[(N)] $\Theta\vdashG(\ev A\rho)\sigma=\ev A{(\rho\sigma)}$ for
  $\vdashG\sigma:\Theta\to\Delta$, where $(\rho\sigma)(i)=\rho(i)\sigma$;
\item[(E)] $\Delta\vdashG\ev A{\rho_1}=\ev A{\rho_2}$ for related $\rho_1,\rho_2$;
\item[(S)] $\Delta\vdashG\ev{A[\alpha]}\rho=\ev A{(\ev\alpha\rho)}$, for
  substitutions $\alpha$ whose components satisfy (P), (N), (E).
\end{itemize}
Terms have the same four properties. For equality judgements the statement is
\emph{relational}: related environments give equal values. Implicit
substitution (S) and explicit substitution (N) are treated separately: the
naturality of the GAT is used only in (N), never in (S).

\begin{lemma}{18}
Every derivation of $T_{\mathsf{Tarski}}$ satisfies (P), (N), (E) and (S),
or the relational statement for an equality.
In particular $\Gamma\vdashT J$ implies $\lceil\Gamma\rceil\vdashG\lceil J\rceil$,
and $\lceil A[\alpha]\rceil$ is $\lceil A\rceil$ evaluated in $\lceil\alpha\rceil$.
\end{lemma}

\begin{proof}[Proof sketch]
By one induction on derivations. The motive also carries the four properties
of the types of the context. Every rule has a premise in the context of its
conclusion, so this information is always available.

Three relations between environments carry the proof. Each is an instance
of Ehrhard's equations:
\begin{align*}
(\up\rho)\,\lift(\sigma)&\ \sim\ \up(\rho\sigma) &&\text{by } \q\lift(\sigma)=\q,\ \p\lift(\sigma)=\sigma\p
 &&\text{(binders in (N))}\\
(\up\rho)\,\sect(a)&\ \sim\ (\ev a\rho,\rho) &&\text{by } \q\sect(a)=a,\ \p\sect(a)=\id
 &&\text{(application, $\beta$)}\\
\ev{\alpha^+}{(\up\rho)}&\ \sim\ \up(\ev\alpha\rho) &&\text{by (N) for } \alpha(i) \text{ and } \sigma=\p
 &&\text{(binders in (S))}
\end{align*}
Application uses $B\,\sect(a)=\ev B{(\ev a\rho,\rho)}=\ev{B[a]}\rho$, by (N),
(E) and then (S) for $B$. Its naturality uses
$\sect(a)\sigma=\lift(\sigma)\sect(a\sigma)$, and $\eta$ uses
$\lift(\p)\sect(\q)=\id$. For a term, (N) and (S) are stated at the type
$\ev A{(\rho\sigma)}$ (resp.\ $\ev A{(\ev\alpha\rho)}$) rather than
$(\ev A\rho)\sigma$. With that choice the conversion rule needs no naturality
of the type.
\end{proof}

\section{Lemma 19 and Theorem 20}

\begin{lemma}{19}
By induction on derivations of $T_{\mathsf E}$, using Theorem~16 and then
Lemma~18 on the erased derivations:
\begin{itemize}
\item $\vdashG\Gamma$ implies $\vdashG\Gamma=\lceil\er\Gamma\rceil$, and $\idv$ is typed over $\Gamma$ for $\er\Gamma$;
\item $\Gamma\vdashG A$ implies $\Gamma\vdashG A=\lceil\er A\rceil$, and
$\Gamma\vdashG a:A$ implies $\Gamma\vdashG a=\lceil\er a\rceil:A$;
\item $\vdashG\sigma:\Delta\to\Gamma$ implies that $\idv\sigma$ and $\ev{\er\sigma}\idv$ are related over $\Delta$.
\end{itemize}
\end{lemma}
The substitution clause is the pointwise form of $\sigma=\lceil\er\sigma\rceil$.
Its cases are the generators. For $\id$ and $\p$ it holds by definition.
The cases $\lift$ and $\sect$ are the first two relations of Lemma~18.
Composition uses (S) for the components of $\er\sigma$.

\begin{theorem}{20}
If $\Gamma\vdashT J$ then $\lceil\Gamma\rceil\vdashG\lceil J\rceil$, with
$\er{\lceil\Gamma\rceil}=\Gamma$ and $\er{\lceil J\rceil}=J$. Two judgements
of $T_{\mathsf E}$ over the same context with the same erasure are convertible,
and every context is convertible to $\lceil\er\Gamma\rceil$.
\end{theorem}

\section{The formalisation}

The files are in \seg{UNIT/} and are checked together by \seg{GATCheck.seg}.
They are also part of \seg{All.seg}.
\begin{center}\small
\begin{tabular}{ll}
\seg{TSyntax}, \seg{TTyping}, \seg{TSubst}, \seg{TSubCong}, \seg{TSubLib} & $T_{\mathsf{Tarski}}$, Lemma 15\\
\seg{GAT}, \seg{GSound} & $T_{\mathsf E}$, erasure, Theorem 16\\
\seg{Ev} & evaluation, Lemma 17\\
\seg{L18Base} \dots\ \seg{L18BE}, \seg{L18} & Lemma 18\\
\seg{L19Base}, \seg{L19}, \seg{Thm20} & Lemma 19, Theorem 20\\
\seg{TPresup}, \seg{TLean} & presuppositions of $T_{\mathsf{Tarski}}$, lean rules\\
\seg{GPresup}, \seg{GLean} & presuppositions of $T_{\mathsf E}$, lean rules\\[3pt]
\seg{LTypes} \dots\ \seg{LSoundF}, \seg{TSem}, \seg{TSound} & levelled filter model, soundness\\
\seg{TInv} & inversion (Theorems 9, 10 for $T_{\mathsf{Tarski}}$)\\
\seg{RTyping}, \seg{ErR} & $T_{\mathsf{Russell}}$, erasure, Lemma 21\\
\seg{TGen} & generation lemmas\\
\seg{TUniqBase}, \seg{TUniqComb}, \seg{TUniq} & Lemmas 25, 26, Corollary 27\\
\seg{TSection} & Theorem 28\\
\seg{RGAT} & $T_{\mathsf{Russell}}$ and $T_{\mathsf E}$ (Section~\ref{sec:compose})\\
\seg{AGAT}, \seg{AForget}, \seg{AUniq}, \seg{AAnnot}, \seg{ARGAT} & context annotations (Section~\ref{sec:annot})\\
\seg{EhrhardCwF}, \seg{EhrhardTEModel} & model interface\\
\seg{TEDenote}, \seg{TEInterpSound} & interpretation and soundness of $T_{\mathsf E}$\\
\seg{RSem}, \seg{RSemV} & Russell semantics (Section~\ref{sec:rsem})\\
\seg{russell-posets/PosetTEModel} & univalent model, one universe (Section~\ref{sec:rsem})\\
\seg{FGAT}, \seg{FForget}, \seg{FPresup}, \seg{FGen} & full syntax, erasure and generation\\
\seg{FCoherence}, \seg{FAnnot}, \seg{FRGAT} & full annotations (Section~\ref{sec:fullannot})
\end{tabular}
\end{center}
Lifts $t^n_m(a)$ are strict ($m<n$), as the cumulativity of $T_{\mathsf{Russell}}$ in Appendix~A.3: with $t^n_n$ and no equation $t^n_n(a)=a$, the terms $t^0_0(a)$ and $a$ have the same Russell erasure without being convertible, and the equivalence with $T_{\mathsf{Russell}}$ fails.

The syntactic proofs use no postulates and no axioms. Some rules carry extra premises,
all of them presuppositions.
In $T_{\mathsf{Tarski}}$:
\begin{itemize}
\item the typing of application carries $\Gamma\vdash A$;
\item the binder congruences carry $\Gamma\vdash A$ (or $\Gamma\vdash a:\U_n$);
\item $\beta$ and $\eta$ carry $\Gamma\vdash A$ and $\Gamma.A\vdash B$;
\item the congruence of application carries $\Gamma.A\vdash B$ and $\Gamma\vdash a:A$;
\item symmetry of term equality carries $\Gamma\vdash A$.
\end{itemize}
In $T_{\mathsf E}$ the same, and in addition $a\sigma$ carries $\Gamma\vdashG A$
(as in Cartmell's presentation, where $\mathrm{Tm}(\Gamma,A)$ presupposes
$A:\mathrm{Ty}(\Gamma)$), the congruences of $A\sigma$, $a\sigma$, $\lift_A$
carry the typings of their left-hand sides, and the congruence of $\sect_A$ carries
$\Gamma\vdashG A$, $\Gamma\vdashG a:A$, $\Gamma\vdashG a':A$.

\section{Presuppositions}\label{sec:presup}
These premises are admissible, so the theorems hold for the lean systems too.

\begin{lemma}[Lemmas 22--24 of the paper]
In $T_{\mathsf{Tarski}}$: $\Gamma\vdash a:A$ implies $\Gamma\vdash A$;
$\Gamma\vdash A=B$ implies $\Gamma\vdash A$ and $\Gamma\vdash B$;
$\Gamma\vdash a=b:A$ implies $\Gamma\vdash A$, $\Gamma\vdash a:A$, $\Gamma\vdash b:A$.
\end{lemma}
One induction on derivations. The binder cases use context conversion; application,
$\beta$ and the congruence of application use the substitution lemma
($\Gamma\vdash B[a]=B'[a']$ from $\Gamma.A\vdash B=B'$ and $\Gamma\vdash a=a':A$).

\begin{lemma}
In $T_{\mathsf E}$ every judgement implies its presuppositions: $\vdashG\Gamma$ for
$\Gamma\vdashG J$, $\Gamma\vdashG A$ for $\Gamma\vdashG a:A$, $\vdashG\Delta$ and $\vdashG\Gamma$
for $\vdashG\sigma:\Delta\to\Gamma$, and both sides typed for every equality.
\end{lemma}
Since substitutions are explicit, there is no meta-level substitution lemma:
each case is a short derivation from the rules. For instance
$(A\p)\,\sect_A(a)=A(\p\,\sect_A(a))=A\,\id=A$ types $\q\,\sect_A(a)$, and
$(B\lift_A(\p))\,\sect_{A\p}(\q)=B$ types the right-hand side of $\eta$.
The equations $\lift_A(\id)=\id$, $\lift_A(\sigma\delta)=\lift_A(\sigma)\lift_{A\sigma}(\delta)$,
the congruences of $\lift$ and $\sect$, and the congruence of application need
the conversion rule for substitutions above.

The lean systems ($T_{\mathsf{Tarski}}$ as in Appendix~A.2, $T_{\mathsf E}$ without the
premises listed) are translated into the systems of the formalisation by recovering
the premises from these lemmas. For $\eta$ this also needs inversion of
$\Gamma\vdash\Pi(A,B)$, which holds in $T_{\mathsf{Tarski}}$ because only one rule
concludes it, and in $T_{\mathsf E}$ by induction, going through the rule converting
types along $\vdashG\Gamma=\Delta$. One difference with Appendix~A.2 remains: symmetry and
transitivity are two premise-free rules, where A.2 has the single rule
$A=B,\ A=C\Rightarrow B=C$.

\section{Inversion for $T_{\mathsf{Tarski}}$}\label{sec:inv}

Section~5 needs only the following inversion properties. Injectivity of
$\Pi$ is \emph{not} needed, because the syntax of $T_{\mathsf{Russell}}$
is annotated.

\begin{lemma}{(Theorems 9, 10)}
In $T_{\mathsf{Tarski}}$:
$\Gamma\vdashT\U_m=\U_n$ implies $m=n$;
$\Gamma\vdashT\Pi(A,B)=\U_n$ is not derivable;
$\Gamma\vdashT\U^k_m=\U^k_n:\U_k$ implies $m=n$;
$\Gamma\vdashT\U^k_m=\Pi_k(a,b):\U_k$ is not derivable.
\end{lemma}
These follow from soundness in a filter model. The model is the domain of
Carneiro, Coquand et al.\ (``Definitional inversion, without normalisation'')
with a hierarchy of codes $\mathsf{univ}\,n$. We interpret
$\ev{\U_n}{}=\ev{\U^n_m}{}$ as the filter of $\mathsf{univ}\,n$
(resp.\ $\mathsf{univ}\,m$); $\T_n$ and $t^n_m$ are interpreted as the
identity, and $\Pi_n$ as $\Pi$. Consistency forbids two different levels,
and also a level together with a $\Pi$-type. We evaluate in the bottom
environment, which fits every well-formed context.

\section{$T_{\mathsf{Russell}}$ and erasure (Lemma 21)}

The raw syntax of $T_{\mathsf{Russell}}$ is the fragment
$v_i,\ \Pi(A,B),\ \lambda(A,B,b),\ \app(A,B,c,a),\ \U_n$ of the Tarski
syntax, so substitution and its laws are shared. The rules are those of
Appendix~A.3, with strict cumulativity: $a:\U_m$ and $m<n$ give $a:\U_n$.
Erasure is
\[
\erR{\T_n(a)}=\erR{t^n_m(a)}=\erR a,\qquad
\erR{\Pi_n(a,b)}=\Pi(\erR a,\erR b),\qquad
\erR{\U^n_m}=\U_m,
\]
and homomorphic otherwise. It commutes with renaming and substitution.

\begin{lemma}{21}
If $\Gamma\vdashT J$ then $\erR\Gamma\vdashR\erR J$.
\end{lemma}
The proof is one induction. As in Section~4, the extra premises of
$T_{\mathsf{Tarski}}$ are mirrored in $T_{\mathsf{Russell}}$, so that the
erasure of each Tarski rule is an instance of a Russell rule. The six Tarski
equations erase to reflexivity.

\section{Uniqueness (Lemmas 25, 26 and Corollary 27)}

For $k\le n$ write $\hat t^n_k(v)$ for $v$ if $k=n$ and for $t^n_k(v)$ if
$k<n$. This is a meta-level operation, since the lifts of the syntax are
strict.

\begin{lemma}{25}
In $T_{\mathsf{Tarski}}$:
\begin{enumerate}
\item $\Gamma\vdashT A$, $\Gamma\vdashT B$ and $\erR A=\erR B$ imply $\Gamma\vdashT A=B$;
\item $\Gamma\vdashT B$, $\Gamma\vdashT a:\U_n$ and $\erR B=\erR a$ imply $\Gamma\vdashT B=\T_n(a)$;
\item $\Gamma\vdashT u_0:A_0$, $\Gamma\vdashT u_1:A_1$ and $\erR{u_0}=\erR{u_1}$
imply one of two things:
either $\Gamma\vdashT A_0=A_1$ and $\Gamma\vdashT u_0=u_1:A_0$;
or $\Gamma\vdashT A_i=\U_{n_i}$ and there is one code $\Gamma\vdashT v:\U_k$,
with $k\le n_0,n_1$ and $\Gamma\vdashT u_i=\hat t^{n_i}_k(v):\U_{n_i}$ for $i=0,1$.
\end{enumerate}
\end{lemma}
The second outcome has one common code $v$ where the paper has two
convertible codes $v_0=v_1:\U_k$. The two forms are equivalent by the
congruence of lifts.

\begin{proof}[Proof sketch]
The three statements are proved together, by induction on
$|u_0|+|u_1|$ (respectively $|A|+|B|$ and $|B|+|a|$), where $|t|$ is the size
of the raw term. This replaces the paper's nested inductions on $u_0$ and
$u_1$, which are needed because an erasure such as $\erR{u_1}=\Pi(A,B)$
does not fix the shape of $u_1$. Derivations are not used as the induction
measure. Instead, \emph{generation lemmas} read off the premises of the
last non-conversion rule, together with a conversion from its type to the
given type. So there is no recursion through the conversion rule.
Each recursive call is on strict subterms:
\begin{itemize}
\item 25.1 on $(X,\T_n(b))$ calls 25.2 on $(X,b)$;
\item 25.2 on $(\T_l(a_1),a_2)$ calls 25.3 on $(a_1,a_2)$; on $(B,t^n_m(a))$ it calls 25.2 on $(B,a)$;
\item 25.2 on $(\Pi(A,B),\Pi_l(a,b))$ calls 25.2 on $(A,a)$ and on $(B,b)$;
\item 25.3 on $(t^n_m(a),u_1)$ calls 25.3 on $(a,u_1)$.
\end{itemize}
The lift case on the right follows from the left one by symmetry. The cases
where the shapes do not match (for instance a $\lambda$ and a lift of a
code) are refuted by the inversion lemma: they would give
$\Gamma\vdashT\Pi(C,B)=\U_n$. Two codes $\U^n_m$ and $\U^{n'}_m$ have the
common code $\U^{m+1}_m$. For two $\Pi$-codes
$\Pi_n(a,b)$ and $\Pi_{n'}(a',b')$, the induction hypothesis gives common
codes $v_a:\U_{k_a}$ and $v_b:\U_{k_b}$. Their common code is
$\Pi_K(\hat t^K_{k_a}(v_a),\hat t^K_{k_b}(v_b))$ with $K=\max(k_a,k_b)$,
using $t^p_n(\Pi_n(a,b))=\Pi_p(t^p_n(a),t^p_n(b))$. This avoids the
trichotomy on $k_a,k_b$.
\end{proof}

\begin{lemma}{26}
In the second outcome of 25.3, if moreover $\Gamma\vdashT A_0=A_1$, then
$\Gamma\vdashT u_0=u_1:A_0$.
\end{lemma}
The proof is as in the paper: $\U_{n_0}=\U_{n_1}$ gives $n_0=n_1$ by
injectivity of universes.

\begin{lemma}{(Corollary 27)}
If $\Gamma\vdashT u_0:A_0$, $\Gamma\vdashT u_1:A_1$, $\erR{u_0}=\erR{u_1}$ and
$\erR{A_0}=\erR{A_1}$, then $\Gamma\vdashT A_0=A_1$ and $\Gamma\vdashT u_0=u_1:A_0$.
\end{lemma}

\section{Section of erasure (Theorem 28)}\label{sec:section}

\begin{theorem}{28}
If $\Gamma\vdashR J$, then there is a judgement $\Gamma'\vdashT J'$ with
$\erR{\Gamma'}=\Gamma$ and $\erR{J'}=J$. It is unique up to conversion. Two
well-formed contexts with the same erasure are convertible. Two lifts of
the same judgement are convertible componentwise; for terms this means
$\Gamma'\vdashT A'=A''$ and $\Gamma'\vdashT a'=a'':A'$.
\end{theorem}

\begin{proof}[Proof sketch]
For existence we prove a stronger statement, by induction on the Russell
derivation of $\Gamma\vdashR J$:
\begin{enumerate}
\item[(a)] there is $\Gamma'$ with $\vdashT\Gamma'$ and $\erR{\Gamma'}=\Gamma$;
\item[(b)] for \emph{every} $\Delta$ whose types are well formed and with
$\erR\Delta=\Gamma$, there is a lift $J'$ of $J$ with $\Delta\vdashT J'$
and $\erR{J'}=J$.
\end{enumerate}
The theorem is (b) applied to the context given by (a).

The paper applies the induction hypothesis to two premises and obtains
lifts in two contexts $\Gamma'$ and $\Gamma''$. One premise must then be
moved into the other context, and the Agda prototype has a separate stage
for this. Fixing the target context $\Delta$ in advance removes that stage.
In the case of $\Pi$, we lift $A$ to $A'$ in $\Delta$ and then apply (b)
for $B$ to $\Delta.A'$, whose erasure is $\Gamma.A$.

What remains is to align types. A premise gives $\Delta\vdashT a':T$ with
$\erR T=A$, and another gives $\Delta\vdashT A'$ with $\erR{A'}=A$. Then
Lemma~25.1 gives $\Delta\vdashT T=A'$, and we apply conversion. When the
type is a universe, $A'$ is $\U_n$. This is how a lifted code is put at
$\U_n$ before applying $\T_n$, $\Pi_n$ or $t^n_m$; for
$\Pi_n(A,B)$ the context is extended by $\T_n(a')$.

Transitivity of term equality needs more. The lifts $a'=b':T$ and
$b''=c'':T'$ agree only up to erasure in the middle term, and
Corollary~27 gives $T=T'$ and $b'=b'':T$. The variable cases split $\Delta$
and weaken; they use $\erR{\mathrm{wk}\,X}=\mathrm{wk}\,\erR X$. The cases
$\beta$, $\eta$ and application use the commutation of erasure with
substitution.

For uniqueness, contexts go by induction: the types at each position are
convertible by Lemma~25.1, after moving one of them along the conversion
of the prefixes (Lemma~24). A lift in $\Gamma''$ is then moved to
$\Gamma'$ by context conversion, and Lemma~25.1 or Corollary~27 applies.
\end{proof}

Together with Lemma~21, this shows that $T_{\mathsf{Tarski}}$ and
$T_{\mathsf{Russell}}$ are equivalent.

\section{$T_{\mathsf{Russell}}$ and $T_{\mathsf E}$}\label{sec:compose}

Write $\|X\| = \erR{\er X}$ for the double erasure from $T_{\mathsf E}$ to
$T_{\mathsf{Russell}}$ (file \seg{RGAT}).

\begin{theorem}{(Russell and $T_{\mathsf E}$)}
\begin{enumerate}
\item If $\vdashG\Gamma$ (resp.\ $\Gamma\vdashG J$) then
$\vdashR\|\Gamma\|$ (resp.\ $\|\Gamma\|\vdashR\|J\|$).
\item If $\vdashR\Gamma$ (resp.\ $\Gamma\vdashR J$) then there are
$\Gamma'$, $J'$ with $\vdashG\Gamma'$ (resp.\ $\Gamma'\vdashG J'$),
$\|\Gamma'\|=\Gamma$ and $\|J'\|=J$.
\item If $\Gamma\vdashG A_1$, $\Gamma\vdashG A_2$ and $\|A_1\|=\|A_2\|$,
then $\Gamma\vdashG A_1=A_2$. If $\Gamma\vdashG a_i:A_i$ with
$\|a_1\|=\|a_2\|$ and $\|A_1\|=\|A_2\|$, then $\Gamma\vdashG A_1=A_2$ and
$\Gamma\vdashG a_1=a_2:A_1$. For equality judgements over $\Gamma$ with the
same double erasure, the corresponding sides are convertible.
\item If $\vdashG\Gamma_1$, $\vdashG\Gamma_2$ and $\|\Gamma_1\|=\|\Gamma_2\|$,
then $\vdashG\Gamma_1=\Gamma_2$.
\end{enumerate}
\end{theorem}
\begin{proof}
(1) is Theorem~16 followed by Lemma~21. (2) is Theorem~28 followed by
Theorem~20: take $\Gamma'=\lceil\Delta\rceil$ and $J'=\lceil J''\rceil$ for the
Tarski lift $\Delta\vdashT J''$. For (3), erase to $T_{\mathsf{Tarski}}$ over
$\er\Gamma$. Lemma~25.1 (resp.\ Corollary~27) gives
$\er\Gamma\vdashT\er{A_1}=\er{A_2}$ (resp.\ also
$\er\Gamma\vdashT\er{a_1}=\er{a_2}:\er{A_1}$). Theorem~20 lifts this to
$\lceil\er\Gamma\rceil$, and we move it to $\Gamma$ along
$\vdashG\lceil\er\Gamma\rceil=\Gamma$. Lemma~19 gives
$\Gamma\vdashG A_i=\lceil\er{A_i}\rceil$ and
$\Gamma\vdashG a_i=\lceil\er{a_i}\rceil:A_i$, and we conclude by transitivity
and conversion; equality judgements reduce to this through their
presuppositions. For (4), the Tarski contexts $\er{\Gamma_1}$ and $\er{\Gamma_2}$
are convertible (uniqueness in Theorem~28), and context conversion of
$T_{\mathsf{Tarski}}$ lifts pointwise along $\lceil-\rceil$ by Theorem~20. This gives
$\Gamma_1=\lceil\er{\Gamma_1}\rceil=\lceil\er{\Gamma_2}\rceil=\Gamma_2$, and we
conclude by symmetry and transitivity of context equality.
\end{proof}
Symmetry and transitivity of context equality are not rules of $T_{\mathsf E}$;
both are proved by induction on derivations. For symmetry, in the case of
$\vdashG\Gamma.A=\Delta.B$ the equation $B=A$ is moved from $\Gamma$ to
$\Delta$. For transitivity, the second derivation is inverted: only the rules
for the empty context, reflexivity and extension conclude a context equality.

\section{Semantics of $T_{\mathsf{Russell}}$}\label{sec:rsem}

The checked versions of the three theories have levels $0,\ldots,N$.
At $N=0$ there is one universe. The preceding sections give the rule
schemes without a global bound. Each derivation uses finitely many levels,
so it admits a finite bound. The formal comparison with the unbounded
presentation is open.

A $\seg{TEModel}[N]$ is a category with families in Ehrhard's presentation
whose contexts, substitutions, types and terms form sets. It has dependent
products, given by a term bijection natural under substitution, and a
cumulative Tarski hierarchy through level $N$. The hierarchy has product
and universe codes, decoding and lifts satisfying the $T_{\mathsf E}$ equations.

The partial interpretation of raw $T_{\mathsf E}$ expressions has
propositional domains, so denotation is functional. A context denotes a
semantic telescope that records each successive type, supplying the last
extension for $\p$ and $\q$. Over a supplied telescope, a term denotes a
type and an element of it. A substitution returns a target telescope
and a model substitution. In $\sigma\delta$, the target of $\delta$
supplies the source for $\sigma$.

Soundness is proved by induction on $T_{\mathsf E}$ derivations.
Well-formed contexts have denotations, and convertible contexts denote the
same telescope. At every denotation of the stated context, types and terms
have values and the two sides of an equation share a value. For substitutions
the computed target must equal the supplied denotation of the target context.

\begin{theorem}{(Russell semantics)}
For every $N$ and every $\seg{TEModel}[N]$:
\begin{enumerate}\setlength{\itemsep}{0pt}
\item Each bounded $T_{\mathsf{Russell}}$ judgement has a denotation
through a $T_{\mathsf E}$ lift under double erasure.
\item Any two lifts of that judgement have the same semantic value.
For types and terms the value includes the telescope, and for terms
it also includes the semantic type. An equality judgement has the common
value of its two sides.
\item Russell-convertible types and Russell-convertible terms have
equal denotations, including when their denotations were obtained from
independent typing derivations.
\end{enumerate}
\end{theorem}
\begin{proof}
For (1), apply the lifting theorem of Section~\ref{sec:compose} and
$T_{\mathsf E}$ soundness. For (2), equality of the lifted contexts
aligns the telescopes. Convert the second typing judgement to the first
context and apply the uniqueness clauses of Section~\ref{sec:compose}.
Soundness gives a common value in the aligned telescope, and functionality
identifies both values with it. For equality judgements, compare the value
of either side. For (3), lift and interpret the Russell equality. By (2)
its common value agrees with the values obtained from any lifts of the
individual typing judgements.
\end{proof}

Gylterud's tower $V_0,\ldots,V_N$ supplies the codes and lifts of a
$\seg{TEModel}[N]$. Model assembly from a supplied tower and the Russell
semantics theorem are uniform in $N$. Unit supplies the tower as a
schema for closed $N$, since each tower uses a statically specified
sequence of Seg meta-levels. The instances for $N=0,1,2$ are checked in
\seg{RSemV}.

The tower is built from W-types in Unit. Univalence is used once, to show
that iterative sets form a set, and it is a parameter of the model. In
this sense the model is a set-theoretic semantics in the style of Aczel
and Werner, with $V_n$ playing the role of an inaccessible. The rule
$a:\U_m\Rightarrow a:\U_n$ of $T_{\mathsf{Russell}}$ is interpreted by the
lift $V_m\to V_n$ of codes. This lift is not the identity, but decoding
does not see it: $\T_n(t^n_m(a))$ and $\T_m(a)$ denote the same type. This
is the situation where the detour through $T_{\mathsf E}$ is needed: a
direct interpretation of Russell derivations would have to decide where
the lifts go, and prove that choice coherent. Theorem~28 and uniqueness
do both once, syntactically.

\paragraph{Other models.} The theorem holds for every
$\seg{TEModel}[N]$, that is, every category with families in Ehrhard's
form whose four sorts are sets, with $\Pi$ and a cumulative Tarski
hierarchy satisfying the equations of $T_{\mathsf E}$. It is not specific
to iterative sets. Context annotations play no role in it, since
$T_{\mathsf E}$ derivations carry their contexts. Annotations matter for a
different statement, namely that the syntax is the initial model of a
generalised algebraic theory, so that the interpretation in a model is the
\emph{unique} morphism of categories with families. That statement is not
formalised here; Section~\ref{sec:annot} relates the annotated syntax to
$T_{\mathsf E}$.

\paragraph{A univalent model with one universe.} For $N=0$ the rules for
$\U^n_m$, $t^n_m$ and cumulativity cannot apply, and a
$\seg{TEModel}[0]$ is an Ehrhard category with families with $\Pi$ and
one universe closed under $\Pi$. The univalent model in
\seg{CANONICITY/unit/sag-model} is already in Ehrhard's form: it has the
category with families of fibrant types over posets (\seg{FibCwF}, with
sets as sorts), $\Pi$ (\seg{FPi}) and one universe $\U$ with decoding
(\seg{CwFU}, instance \seg{FibUniv}). Here $\U(X)$ consists of all small
fibrant types over $yX$, so every small fibrant type $A$ over $\Gamma$ has
a code $\code{A}:\Gamma\to\U$, $\rho\mapsto A\hat\rho$ with its structure
moved along $\hat\rho$. Following Huber (thesis, Definition~1.1.4), the
strict equation $\T(\code{A})=A$, as fibrant types, makes codes for the
type formers unnecessary: $\Pi_0(a,b):=\code{\Pi(\T a,\T b)}$ satisfies
$\T(\Pi_0(a,b))=\Pi(\T a,\T b)$.

\begin{theorem}{(univalent model)}
The universe of the poset model is closed under $\Pi$ with strict
decoding, natural in substitutions. Hence the model is a
$\seg{TEModel}[0]$, and the Russell semantics theorem holds in it.
\end{theorem}
\begin{proof}
On the underlying presheaves, $\T(\code{A})=A$ is the code-map lemma of
the model, up to the identity law of composition, which suffices since
types form a set. For the fibration structures, restrict along a Yoneda
point $\hat\rho$: the code map restricted is the Yoneda code of
$\code{A}\rho$, and the lemma \emph{El along a Yoneda code}
(\seg{CPosEl5}) gives that the fibration of $\mathsf{El}$ moved along it is
that of $A$ moved along $\hat\rho$. Two fibration structures that agree
after moving along every Yoneda point are equal (\seg{FibLocal}): a
relative contractibility structure is determined by its values, the value
of a moved structure at the generic point is the original value at
$\rho$, and since the interval is tiny on posets a composition problem
at $X$ is reached from the generic point of $X\times[1]$. This gives
$\T(\code{A})=A$ (\seg{FibName}, \seg{FibNameDec}), and also
$\code{\T c}=c$ (\seg{FibUnivInj}). The $\Pi$-type of small types,
built at the level of the universe, equals the $\Pi$-type of the model
through the lift of iterative sets (\seg{FibPiSmall}). Naturality of
codes, $\code{A}\sigma=\code{A\sigma}$, gives naturality of $\Pi_0$
(\seg{FibUnivPi}). At $N=0$ the remaining fields of
$\seg{TEModel}[0]$ concern pairs of levels $m<n\le 0$ and are vacuous;
contexts form a set because a presheaf context is determined by its codes
and action (\seg{PosetContextSet}). The instance is
\seg{PosetTEModel}.
\end{proof}

The files of the classifier are in \seg{sag-model}. The instance and the
Russell semantics at $N=0$ are in \seg{UNIT/russell-posets}, checked
together with \seg{core} and \seg{sag-model} by its \seg{check.py}
(15288 declarations, no failure) with the indexed checker. That checker
now resolves \texttt{opaque} and \texttt{transparent} locks as the main
checker does, so both accept the same sources. A closed example,
$\lambda(A:\U_0).\,\Pi(x:A).\,A$ at type $\Pi(A:\U_0).\,\U_0$, checks
that a direct derivation and one ending in a conversion have the same
denotation (\seg{PosetRSemCheck}). The model is relative to function
extensionality and univalence, which are parameters of the poset model.
The syntax has no identity types, so univalence is not visible in
$T_{\mathsf{Russell}}$; the result shows that the semantics applies in a
model that is not set-theoretic. Adding identity types to the four
theories is a separate project.

\section{Context annotations}\label{sec:annot}

The theory $T_{\mathsf E}^{a}$ gives context arguments to the six leaves
$(\U_n)_\Gamma$, $(\U^n_m)_\Gamma$, $\id_\Gamma$, $()_\Gamma$,
$\p_{\Gamma,A}$ and $\q_{\Gamma,A}$, as in Section~2.1 of
Coquand--Sterbac. It uses Ehrhard's $\lift_A$ and $\sect_A$ and the same
level bound $N$ as $T_{\mathsf E}$. Each leaf has a congruence rule in its
annotation arguments. For example, $\Gamma=\Delta$ and $\Gamma\vdash A=B$
give $\p_{\Gamma,A}=\p_{\Delta,B}:\Gamma.A\to\Gamma$ and
$\Gamma.A\vdash\q_{\Gamma,A}=\q_{\Delta,B}:A\p_{\Gamma,A}$ in the
annotated theory. The right projection is typed at the left endpoints
by context conversion. The right variable is moved to $\Gamma.A$
and its type is converted using the projection equation.

Write $\mathrm{fg}$ for the map that drops the six leaves' annotation
fields. It preserves derivability and sends the annotation congruences
to reflexivity. The converse takes a derivable bare judgement $J$ to an
annotated judgement $J^a$ with $\mathrm{fg}(J^a)=J$ as raw judgements.
Any well-formed annotated representatives of the contexts of $J$ can
be prescribed, including both endpoints of a substitution or a
context equation.

The lifting proof uses coherence for expressions with identical raw
images under $\mathrm{fg}$. Generation recovers the introduction premises
through context and type conversion. One induction on the common bare
expression compares types, terms, substitutions and well-formed
contexts. In the substitution case, equality of the domains gives both
equality of the codomains and a substitution equation at the left
endpoints. This supplies the context equality needed to compare $A$ and
$A'$ in the case of $A\sigma$ and $A'\sigma'$. The presupposition
induction also supplies symmetry of context equality. The lifting
induction then aligns independently lifted premises using coherence,
and context conversion places the resulting judgement over any
prescribed representatives.

\begin{theorem}{(Russell and $T_{\mathsf E}^{a}$)}
The four conclusions of Section~\ref{sec:compose} hold for
$T_{\mathsf E}^{a}$ with erasure \mbox{$X\mapsto\|\mathrm{fg}(X)\|$}.
\end{theorem}
\begin{proof}
For soundness, forget annotations and apply the bare erasure theorem.
For existence, apply the bare Russell lift and then annotate, using the
raw identity $\mathrm{fg}(J^a)=J$. For uniqueness, apply the bare
uniqueness theorem to the forgetful images of the given annotated
judgements. Lift the resulting equations into the given annotated
contexts. Coherence aligns the lifted endpoints with the original
types and terms. Context equality lifts directly between the two
given annotated contexts. The composition is checked in \seg{ARGAT}.
\end{proof}

\section{Complete sort annotations}\label{sec:fullannot}

In Cartmell's style, $T_{\mathsf E}^{\mathrm{full}}$ keeps the six leaf
annotations and gives the formers the ambient context
$\Pi_\Gamma$, $\lambda_\Gamma$, $\app_\Gamma$, $\T_{n,\Gamma}$,
$\Pi_{n,\Gamma}$ and $t^n_{m,\Gamma}$. For $\delta:\Theta\to\Delta$,
$\sigma:\Delta\to\Gamma$ and $\Gamma\vdash a:A$, the other operations
have the following sorts, with separate type and term actions.
\begin{align*}
 A[\sigma]_{\Delta,\Gamma} &:\mathrm{Ty}(\Delta),&
 a[\sigma]_{\Delta,\Gamma,A} &:\mathrm{Tm}(\Delta,A[\sigma]_{\Delta,\Gamma}),\\
 \sigma\circ_{\Theta,\Delta,\Gamma}\delta &:\Theta\to\Gamma,&
 \lift_{\Delta,\Gamma,A}(\sigma) &:\Delta.A[\sigma]_{\Delta,\Gamma}\to\Gamma.A,\\
 \sect_{\Gamma,A}(a)&:\Gamma\to\Gamma.A.
\end{align*}
The bound $N$ and Ehrhard's generators are the same as before.

Each new sort argument has a congruence at the left formation frame.
For $\Gamma\vdash A=A'$, the terms $a[\sigma]_{\Delta,\Gamma,A}$
and $a[\sigma]_{\Delta,\Gamma,A'}$ are equal at $A[\sigma]_{\Delta,\Gamma}$.
The map $\mathrm{ff}$ drops the new fields, sends both actions to
the single action in $T_{\mathsf E}^{a}$, and erases the new congruences
to reflexivity.

\begin{theorem}{(Full annotation)}
The map $\mathrm{ff}$ preserves derivability. Each derivable
$T_{\mathsf E}^{a}$ judgement has a full lift with exactly that raw image,
at any prescribed well-formed representatives of its contexts.
Full expressions with the same raw image are convertible at compatible
sorts. For substitutions, equality of the domains also gives equality
of the codomains.
\end{theorem}
\begin{proof}
Generation recovers the introduction frame through conversion.
Coherence is by induction on the common raw image. For composition,
the right operands identify the middle contexts. For term action,
the substitutions identify the operand contexts, and comparison of
the operands identifies their stored types even when those types
have different raw images. The lifting induction uses coherence to
align its premises and context conversion to prescribe its endpoints.
\end{proof}

By Section~\ref{sec:annot}, the four conclusions of
Section~\ref{sec:compose} hold for $T_{\mathsf E}^{\mathrm{full}}$
with erasure \mbox{$X\mapsto\|\mathrm{fg}(\mathrm{ff}(X))\|$}.
For uniqueness, lift the $T_{\mathsf E}^{a}$ equations and align their
sides by coherence. The composition is checked in \seg{FRGAT}.

\end{document}
